\section{Local replacement of nearby support diameters}\label{sec:spread}

In this section we prove Theorem~\ref{thm:spread}. Dividing all summands by
the variance-weighted mean diameter $h_0$, we set
\[
 q_i=1-p_i,\qquad r_i=p_i^2+q_i^2,\qquad
 R_0=\sum_{i=1}^n a_i r_i,\qquad\mu=\sum_{i=1}^n p_i,
 \qquad\ell=\sum_{i=1}^n a_i|e_i|\le\sqrt A\,\sigma,
\]
where $\sum_{i=1}^n a_i e_i=0$ and $A/2\le R_0\le A$. Writing
$K=\sum_{i=1}^n B_i$ and $T=\sum_{i=1}^n e_i(B_i-p_i)$, the centered sum is
$S=K-\mu+T$ with variance $A+\sigma^2$. Set $c_k=\Pp(K=k)$ and
$N_j(k)=\E[T^j\ind_{\{K=k\}}]$.

\subsection{Probability-free Poisson-binomial and weighted-moment bounds}

Uniformly over all probabilities $p_i\in(0,1)$ and all $k\in\Z$, as
$A\to\infty$,
\begin{align}
 c_k&=\frac{\phi(x_k)}{\sqrt A}+O(A^{-1}),\notag\\
 F_K(k)-\Phi(x_k)
   &=\frac{\phi(x_k)}{\sqrt A}
      \left\{\frac12+\frac{\gamma_3}{6}(1-x_k^2)\right\}
           +O(A^{-1}),                              \label{eq:pb-grid}\\
 x_k&=\frac{k-\mu}{\sqrt A},\qquad
 \gamma_3=\frac1A\sum_{i=1}^n a_i(1-2p_i)\in[-1,1],\notag
\end{align}
and the unconditioned weighted moments $N_j(k)=\E[T^j\ind_{\{K=k\}}]$ satisfy
\begin{align}
 |N_1(k)|&\le C\ell A^{-3/2},\notag\\
 |N_2(k)-\sigma^2c_k|&\le \frac{C\sigma^2}{A},\notag\\
 0\le N_4(k)&\le \frac{C(\sigma^4+\eta^2\sigma^2)}{\sqrt A}.
                                                    \label{eq:weighted-coefficients}
\end{align}
To prove \eqref{eq:weighted-coefficients} without dividing by individual
characteristic functions, write $f_i(t)=q_i+p_i e^{\ii t}$, $s=e^{\ii t}-1$,
$\Psi=\prod_{i=1}^nf_i$, and $\Psi_i,\Psi_{ij}$ for the products with one or
two factors omitted. Because $\sum_{i=1}^n a_ie_i=0$, the Fourier transform
of $N_1$ is $\widehat N_1(t)=s\sum_i a_ie_i\Psi_i=-s^2\sum_i a_ip_ie_i\Psi_i$,
whose $L^1([-\pi,\pi])$ norm is $O(\ell A^{-3/2})$. Similarly, using
$\Psi_{ij}=\Psi+(1-f_i)\Psi_i+(1-f_j)\Psi_j+(1-f_i)(1-f_j)\Psi_{ij}$ and
$u_i=a_ie_i$ ($\sum_i u_i=0$) bounds $\|\widehat N_2-\sigma^2\Psi\|_1=O(\sigma^2/A)$,
and expanding $T^4$ into index partitions $(4),(3,1),(2,2),(2,1,1),(1,1,1,1)$
bounds $\|\widehat N_4\|_1=O((\sigma^4+\eta^2\sigma^2)/\sqrt A)$.

\subsection{Threshold localization prior to conditioning}

At any threshold $y=k-\mu+\theta$ with $k=\operatorname{round}(y+\mu)$ and
$|\theta|\le1/2$, decomposing $\{S\le y\}$ by the value of $K$ and summing
$N_4(j)/(|j-k|-1/2)^4$ over $j\ne k$ gives
\begin{equation}\label{eq:crossing-decomposition}
 F_h(y)=F_K(k)-c_k\Pp(T>\theta\mid K=k)+E_{k,\theta},
 \qquad |E_{k,\theta}|\le
             \frac{C(\sigma^4+\eta^2\sigma^2)}{\sqrt A},
\end{equation}
while the third-moment sum $R_h$ satisfies
\begin{equation}\label{eq:resource-change}
 R_h-R_0=3\sum_{i=1}^n a_i r_i e_i+3\sum_{i=1}^n a_i r_i e_i^2
                                     +\sum_{i=1}^n a_i r_i e_i^3,
 \qquad |R_h-R_0|\le C(\ell+\sigma^2+\eta\sigma^2).
\end{equation}
Thus any threshold achieving $\mathcal C_h(y)\ge c_0$ has bounded $|x_k|$,
so $c_k\ge c_{c_0}/\sqrt A$, and dividing \eqref{eq:weighted-coefficients}
by $c_k$ yields
\begin{align}
 |\E(T\mid K=k)|&\le \frac{C\ell}{A},\notag\\
 \E(T^2\mid K=k)&=\sigma^2\bigl\{1+O(A^{-1/2})\bigr\},\notag\\
 \E(T^4\mid K=k)&\le C(\sigma^4+\eta^2\sigma^2).
                                                    \label{eq:conditional-bulk}
\end{align}

\subsection{Conditional $L^1$ lower bound via rejective-sampling matchings}

We next prove that at any localized count $k$,
\begin{equation}\label{eq:conditional-l1}
                   \E(|T|\mid K=k)\ge \frac{c\ell}{\sqrt A}.
\end{equation}
The conditional law of $(B_1,\ldots,B_n)$ given $K=k$ is a rejective
sampling law \citep[Section~2]{BoistardLopuhaaRuizGazen2012}; by
\eqref{eq:pb-grid}, for any set $F$ of at most three indices,
$\Pp(B_F=b_F\mid K=k)=\Pp(B_F=b_F)\{1+O(A^{-1/2})\}$. Let $B,B'$ be two
independent draws from the conditional law given $K=k$, and let
$C=\{i:B_i\ne B_i'\}$, which has even cardinality $m=|C|$. Conditional on
$C$ and on the common values outside $C$, the signs $B_i-B_i'\in\{-1,1\}$ for
$i\in C$ are uniformly distributed over balanced sign vectors (since the
odds product $\prod_{i\in C}(p_i/q_i)$ is independent of which half of $C$ is
$+1$). Generating a uniform balanced sign vector by choosing a uniform
perfect matching of $C$ and assigning independent Rademacher signs to each
matched pair gives, by Khintchine's inequality on each matching,
\[
 \E(|T-T'|\mid C)
 \ge\sqrt{\frac23}\,
       \frac{\sum_{i<j\in C}|e_i-e_j|}{(m-1)\sqrt m}\qquad(m\ge2).
\]
Since $\Pp(i,j\in C)\ge ca_ia_j$ and $\E(m\mid i,j\in C)\le CA$, Jensen's
inequality for $m^{-3/2}$ and $\sum_ja_j|e_i-e_j|\ge A|e_i|$ yield
\[
 2\E(|T|\mid K=k)\ge\E(|T-T'|\mid K=k)
 \ge cA^{-3/2}\sum_{i<j}a_ia_j|e_i-e_j|
 \ge \frac{c\ell}{\sqrt A},
\]
proving \eqref{eq:conditional-l1}.

\subsection{Completion of the comparison and optimality of the rate}

Setting $M_1=\E(|T|\mid K=k)$, $W=\Pp(T>\theta\mid K=k)$, and $L=\theta+W$,
equations \eqref{eq:conditional-bulk}--\eqref{eq:conditional-l1} and Cauchy--Schwarz
give $|\E(T\mid K=k)|\le M_1/4$, $M_1\ge c\sigma^2/(\sigma+\eta)$, and
\begin{equation}\label{eq:relative-phase}
 L\ge c\max\left(\frac\ell{\sqrt A},\,
                      \frac{\sigma^2}{\sigma+\eta}\right)>0,
 \qquad |\theta|+\theta^2+W\le CL,
 \qquad\sigma^2\le C(\sigma+\eta)L.
\end{equation}
Comparing $D_h(y)=F_h(y)-\Phi(y/\sqrt{A+\sigma^2})$ directly with the
unperturbed discrepancy $D_0(k)=F_K(k)-\Phi(x_k)$ via \eqref{eq:crossing-decomposition}
and \eqref{eq:resource-change} shows that every error term is $O((\sigma+\eta)^3L+L/\sqrt A)$,
giving
\[
 \mathcal C_h(y)-\frac{A^{3/2}}{R_0}D_0(k)
 \le\bigl\{-c+C(\sigma+\eta)^3+CA^{-1/2}\bigr\}L
 \le-\kappa\frac{\sigma^2}{\sigma+\eta},
\]
which proves Theorem~\ref{thm:spread} (and when $p_i\in J\subset(0,1)$ for a
compact set $J$, $\eta\le\sigma/\sqrt{\min_{p\in J}p(1-p)}$ gives
$\sigma^2/(\sigma+\eta)\ge c_J\sigma$).

Finally, to see that the rate $\sigma^2/(\sigma+\eta)$ cannot be improved to
$\sigma$ without a lower bound on $p_i$, take a fixed large i.i.d.\ Esseen
group of size $m$ with ratio $>c_0$ attained on the positive closed branch,
and adjoin one rare coordinate $B_{m+1}\sim\Bern(s)$ ($s\to0$) with diameter
$g_s=1-\eta$ ($0<\eta<1/4$) while giving the $m$ group coordinates diameter
$g_c=1+s(1-s)\eta/(ma)$. Because $0<g_s/g_c<1$, the closed CDF at uncentered
threshold $kg_c$ equals $(1-s)F_{\mathrm{group}}(k)+sF_{\mathrm{group}}(k-1)$,
identical to the common-diameter value at $\eta=0$. Since the smooth branch
ratio is analytic in $(s,\eta)$ near $(0,0)$ and is independent of $\eta$ at
$s=0$, its mixed partial derivative is bounded, so
$0<\mathcal C_{\mathrm{com}}(s)-\mathcal C_h\le K_m s\eta$. Choosing $s=\eta^4$
gives $\sigma\sim\sqrt s\,\eta=\eta^3$ and gain $O(\eta^5)=O(\sigma^2/(\sigma+\eta))$.
